% =============================================================================
\section{Main Theoretical Results}\label{sec:theory}
% =============================================================================

\subsection{Basic properties: monotonicity and invariances}

\begin{proposition}[Monotonicity in $\beta$]\label{prop:mono_beta}
For fixed $r$, the mapping $\beta\mapsto \ECL_\beta(r)$ is non-decreasing on $(0,1]$.
\end{proposition}
\begin{proof}
See \cref{app:proofs}.
\end{proof}

\begin{proposition}[Invariance under isometries]\label{prop:isometry}
Assume $d_\cZ(z,z')=\norm{z-z'}_2^2$. If $T:\bbR^d\to\bbR^d$ is a linear isometry (i.e., $T^\top T = I$), then $\ECL_\beta(r)$ is identical for $f$ and $T\circ f$.
\end{proposition}
\begin{proof}
A linear isometry preserves Euclidean distances: $\norm{Tz-Tz'}_2^2 = \norm{z-z'}_2^2$. All influence energies and their cumulative sums are unchanged.
\end{proof}

\begin{remark}[Invariance under general linear maps via whitening]
If $T$ is invertible but not an isometry, Euclidean distances change. Provided $\Sigma := \Cov(f(X))$ is positive definite (i.e., $f(X)$ has full-rank covariance), using Mahalanobis distance $d_\Sigma$ restores invariance under invertible linear transforms, since whitened embeddings $\Sigma^{-1/2}(f(X)-\E[f(X)])$ have identity covariance.
\end{remark}

\subsection{ECL under exponential influence decay}\label{sec:exp_decay}

Effective receptive field theory~\citep{Luo2016EffectiveRF} suggests that influence often decays approximately exponentially (or Gaussian) with distance. We formalize this.

\begin{theorem}[ECL bound under exponential decay]\label{thm:exp_decay}
Assume $\cI(d;r) \le C e^{-\lambda d}$ for constants $C > 0$, $\lambda > 0$, and all $d \ge 0$. Assume positions at distance $d$ have multiplicity $n_d \le 2$ (interior reference). Then
\begin{equation}\label{eq:exp_bound}
\ECL_\beta(r) \;\le\; \frac{1}{\lambda}\log\frac{1}{1-\beta} + \frac{1}{\lambda}\log\frac{2C}{\lambda \cI_{\mathrm{tot}}(r)} + 1.
\end{equation}
In particular, $\ECL_\beta(r) = O\!\left(\lambda^{-1}\log\frac{1}{1-\beta}\right)$ as $\beta \to 1$.
\end{theorem}
\begin{proof}
See \cref{app:proofs}.
\end{proof}

This gives an interpretable formula: ECL scales as $1/\lambda$ (inverse decay rate) times a logarithmic threshold factor. Models with slower decay ($\lambda$ small) have proportionally larger ECL.

\subsection{Architectural locality implies bounded ECL}

\begin{assumption}[Exact locality]\label{ass:exact_local}
There exists $R^\star\in\bbN$ such that for all $x,x'\in\cX$, $x_{W_{R^\star}(r)} = x'_{W_{R^\star}(r)}$ implies $f(x)=f(x')$.
\end{assumption}

\begin{proposition}[ECL upper bound under exact locality]\label{prop:ecl_local}
Under \cref{ass:exact_local}, $\ECL_\beta(r) \le R^\star$ for any $\beta\in(0,1]$.
\end{proposition}
\begin{proof}
If $|i-r|>R^\star$, then perturbing coordinate $i$ leaves $W_{R^\star}(r)$ unchanged, so $\cI(i;r)=0$. Hence $\cI_{\le R^\star}(r)=\cI_{\mathrm{tot}}(r)$.
\end{proof}

\begin{proposition}[Compositional ECL bound]\label{prop:composition}
For an architecture $f = g \circ h$ where $h:\cX \to \cH$ has exact receptive field $R_h$ per position and $g:\cH \to \cZ$ has exact receptive field $R_g$ per position (in the intermediate representation), then $\ECL_\beta(r) \le R_h + R_g$.
\end{proposition}
\begin{proof}
See \cref{app:proofs}.
\end{proof}

\begin{remark}[Strided and dilated stacks]
For a stack of $D$ convolutional layers with kernel sizes $k_l$ and strides $s_l$, the theoretical receptive field is $R^\star = 1 + \sum_{l=1}^D (k_l - 1) \prod_{m<l} s_m$. For a stack of $D$ dilated convolutional layers with constant kernel size $k$ and dilation rates $d_1, \dots, d_D$ (all strides equal to one), the receptive field reduces to $R^\star = 1 + \sum_{l=1}^D (k-1) d_l$. For Basenji's architecture with exponentially increasing dilations, $R^\star = O(k \cdot 2^D)$.
\end{remark}

\subsection{Variance-based sensitivity and Sobol connections}\label{sec:sobol}

We relate perturbation ECL to classical variance decompositions. Let $\varphi:\cZ\to \bbR$ be a scalar summary of embeddings and $Y:=\varphi(f(X))$.

\begin{theorem}[Perturbation influence and Sobol total-effect index, additive case]\label{thm:sobol_equiv}
Assume (i) $f(X) = \sum_{i=1}^L g_i(X_i)$ (additive vector model in $\cZ = \bbR^d$), (ii) coordinates $\{X_i\}$ are independent under $P_X$, (iii) perturbation replaces $X_i$ with an independent copy $X_i' \sim P_i$, and (iv) $d_\cZ = \norm{\cdot}_2^2$. Then:
\begin{enumerate}[leftmargin=2em, label=(\alph*)]
\item (Vector identity.) The single-site influence energy equals twice the trace of the per-coordinate component covariance:
\[
\cI(i;r) \;=\; 2\,\mathrm{tr}\!\big(\Cov(g_i(X_i))\big).
\]
Defining a vector total-effect index $S_i^{T,\mathrm{vec}} := \mathrm{tr}(\Cov(g_i(X_i)))/V_{\mathrm{tot}}^{\mathrm{vec}}$ with $V_{\mathrm{tot}}^{\mathrm{vec}} := \mathrm{tr}(\Cov(f(X)))$, we have $\cI(i;r) = 2\,V_{\mathrm{tot}}^{\mathrm{vec}}\, S_i^{T,\mathrm{vec}}$.
\item (Scalar projection.) For any scalar projection $\varphi:\cZ\to\bbR$ with $Y := \varphi(f(X))$ and $V_{\mathrm{tot}} := \Var(Y)$, the projected single-site influence $\cI^{\varphi}(i;r) := \E[(Y - \varphi(f(X^{(i)})))^2]$ satisfies
\[
\cI^{\varphi}(i;r) \;=\; 2\,\Var(\varphi(g_i(X_i))) \;=\; 2\,V_{\mathrm{tot}}\, S_i^T,
\]
where $S_i^T$ is the classical Sobol total-effect index for coordinate $i$ of $Y$.
\end{enumerate}
\end{theorem}
\begin{proof}
See \cref{app:additive}.
\end{proof}

Under the additive model, ECL based on perturbation influence is thus equivalent to the Sobol total-effect radius. Non-additive models introduce interaction terms; see \cref{sec:interaction_ecl_theory} for quantitative bounds.

\subsection{Tail approximation under weak additivity}\label{sec:interaction_ecl_theory}

\begin{assumption}[Weak interaction/additivity]\label{ass:add}
There exists a decomposition $f(X) = \sum_{i=1}^L g_i(X_i) + \epsilon(X)$, where $\epsilon$ is a small interaction remainder: $\E[\norm{\epsilon(X)}^2] \le \eta^2$.
\end{assumption}

\begin{proposition}[Tail approximation under additivity]\label{prop:tail_add}
Under \cref{ass:add} with independent coordinate resampling and $d_\cZ = \norm{\cdot}_2^2$:
\[
\left|\Delta_{\mathrm{out}}(\ell;r) - \sum_{|i-r|>\ell}\cI(i;r)\right| \;\le\; 4\eta\sqrt{\sum_{|i-r|>\ell}\cI(i;r)} + 4\eta^2.
\]
\end{proposition}
\begin{proof}
See \cref{app:additive}.
\end{proof}

\subsection{Effective receptive field: a linearized derivation}\label{sec:erf}

This subsection records the linear-regime sanity check that the perturbation-variance ECL recovers the squared-impulse-response intuition of effective receptive field theory~\citep{Luo2016EffectiveRF}: in the locally linear case, our influence profile reduces to the squared impulse response, recovering the classical ERF as a special case.

Consider a locally linear model $y_r = \sum_i h_{r-i}\, u_i$ with impulse response $h$.

\begin{proposition}[Influence profile equals squared impulse response]\label{prop:linear_influence}
If perturbation at $i$ replaces $u_i$ by an independent copy with $\E[(u_i-u_i')^2]=\sigma^2$, then $\cI(i;r) = h_{r-i}^2\, \sigma^2$.
\end{proposition}
\begin{proof}
$y_r-y_r^{(i)} = h_{r-i}(u_i-u_i')$. Squaring and taking expectation yields $h_{r-i}^2\sigma^2$.
\end{proof}

\subsection{Model comparison via Blackwell ordering}\label{sec:blackwell}

Fix a radius $\ell$ and define the context-limited embedding $Z_\ell := f(X^{(\mathrm{out},\ell)})$, where $X^{(\mathrm{out},\ell)}$ perturbs outside $W_\ell(r)$. For a downstream decision problem with action space $\cA$, loss $\ell_{\mathrm{dec}}:\cA\times \cY\to\bbR$, and target $Y$, the Bayes risk is
$\cR_f(\ell) := \inf_{\delta:\cZ\to \cA} \E[\ell_{\mathrm{dec}}(\delta(Z_\ell),Y)]$.

\begin{definition}[Blackwell dominance for context experiments]\label{def:blackwell}
Model $f$ \emph{Blackwell-dominates} $g$ at radius $\ell$ if there exists a Markov kernel $K$ such that $g(X^{(\mathrm{out},\ell)}) \stackrel{d}{=} K(f(X^{(\mathrm{out},\ell)}))$.
\end{definition}

\begin{theorem}[Blackwell dominance implies no worse Bayes risk]\label{thm:blackwell}
If $f$ Blackwell-dominates $g$ at radius $\ell$, then $\cR_f(\ell) \le \cR_g(\ell)$ for every decision problem.
\end{theorem}
\begin{proof}
Direct consequence of Blackwell's comparison theorem~\citep{Blackwell1951Comparison}: any decision rule based on $g$ can be simulated from $f$ via $K$.
\end{proof}

\begin{proposition}[ECL ordering does not imply Blackwell dominance]\label{prop:ecl_not_blackwell}
There exist models $f,g$ with $\ECL_\beta(f,r) < \ECL_\beta(g,r)$ for all $\beta$ yet neither Blackwell-dominates the other.
\end{proposition}
\begin{proof}
See \cref{app:proofs} for an explicit counterexample.
\end{proof}

\begin{proposition}[Sufficient condition for risk ordering]\label{prop:risk_ordering}
If $f$ Blackwell-dominates $g$ at every radius $\ell \le \ell_0$, then $\cR_f(\ell) \le \cR_g(\ell)$ for all $\ell \le \ell_0$ and all decision problems. In particular, model $f$ extracts at least as much useful information from any context budget $\ell \le \ell_0$.
\end{proposition}

\subsection{Invariance under sufficient statistics}

\begin{theorem}[ECL invariance under sufficient statistics]\label{thm:sufficient}
Let $g:\cZ\to\cY_{\mathrm{task}}$ be a downstream predictor and $T:\cZ \to \cZ'$ a sufficient statistic for the conditional family $\{P(Y\mid Z = z)\}_{z\in\cZ}$ associated with a fixed task. Then $\ECL^{\mathrm{task}}_\beta(r)$ computed using the composition $g \circ f$ equals $\ECL^{\mathrm{task}}_\beta(r)$ computed using $g \circ T \circ f$.
\end{theorem}
\begin{proof}
Sufficiency implies $P(Y \mid Z) = P(Y \mid T(Z))$, so any decision rule based on $Z$ achieves the same Bayes risk as one based on $T(Z)$. The task-aware influence energies are therefore identical.
\end{proof}

\subsection{Sample-complexity lower bound}

\begin{theorem}[Sample-complexity lower bound for ECL identification]\label{thm:minimax}
Let $\cF_M$ be the class of models with influence variables $U_i$ bounded in $[0,M]$, and consider estimating $\ECL_\beta(r)$ from $n$ i.i.d.\ sequence-perturbation pairs. Then there exist two models $f_0, f_1 \in \cF_M$ with $\ECL_\beta(f_0) \neq \ECL_\beta(f_1)$ and margin $\gamma := \min_{\ell\neq \ECL_\beta(r)}\abs{\cI_{\le \ell}(r) - \beta\,\cI_{\mathrm{tot}}(r)}$ such that any estimator $\hat{\ell}$ achieving
\[
\max_{j\in\{0,1\}}\Pr_{f_j}\!\left(\hat{\ell} \neq \ECL_\beta(f_j)\right) \;\le\; \tfrac{1}{3}
\]
requires
\[
n \;\ge\; c\,\frac{M^2}{\gamma^2}
\]
for an absolute constant $c > 0$. The Monte Carlo estimator described in \cref{sec:algorithms} (\cref{alg:mc_ecl}) attains the matching upper bound on $n$, up to logarithmic factors in $L$ and $1/\delta$, via \cref{cor:sample_complexity}.
\end{theorem}
\begin{proof}
See \cref{app:proofs}. The lower bound follows from a Le Cam two-point construction; the matching upper bound combines \cref{thm:ecl_stability} with Bernstein concentration.
\end{proof}
